

\section{Introduction}

% Plain Language Summarization (PLS) generates summaries of technical documents accessible to non-expert readers. Readability -- the standard measure for PLS -- has long been considered a binary attribute (expert vs. non-expert.) However, readability exists on a spectrum, with different readability levels appropriate for different readers.

Effective scientific summarization must consider multiple factors, including both form and audience~\cite{cohan-etal-2022-overview-first,ScisummNet.v33i01.33017386,BioLaySumm_2024}. When targeting non-expert readers, Plain Language Summarization (PLS) seeks to make complex scientific documents more accessible. This goal is particularly important given the significant public investment in research—approximately 40\% of basic research in the United States is government-funded~\cite{Pece_Anderson_2024}. To promote equitable access to publicly funded research, the U.S. National Institutes of Health now requires all NIH-funded papers to be open access\footnote{\href{https://grants.nih.gov/policy-and-compliance/policy-topics/public-access/nih-public-access-policy-overview}{NIH
 Public Access Policy}}. However, while open access broadens availability, PLS ensures research is understandable to the general public.

%For example, the recent COVID-19 pandemic illustrated the importance of public access to emerging scientific discoveries, as the general population needed to stay up-to-date on the latest understanding of symptoms, how the virus spread, and how to keep themselves and loved ones safe~\cite{wang-etal-2020-cord}. 

% \todoit{while accessibility to all is an important...}
Past work has found that users with higher, but non-expert, familiarity with a topic prefer more technical summaries and found lower complexity summaries less useful~\cite{August2024KnowYA}. However, most prior work in PLS has treated readability as a binary (plain versus technical), aiming for maximum readability, as measured by readability metrics such as Dale-Chall and Flesch-Kincaid Grade Level~\cite{Guo2020AutomatedLL,goldsack2022making,dale1948formula,tanprasert-kauchak-2021-flesch}. To address this gap, we explore fine-grained controls of readability for scientific document summarization using control vectors~\cite{zou2025representationengineeringtopdownapproach, rimsky-etal-2024-steering}, a lightweight controllability method that modifies the model's hidden states. Control vectors are derived from paired examples along a spectrum, capturing the difference in hidden states between the paired examples. During inference, they steer the model towards either end of the spectrum by applying a specified weight. 

 % \todoit{move this sentence up:} 


Control vectors can effectively steer complex attributes like honesty, making them well suited for readability control~\cite{zou2025representationengineeringtopdownapproach, rimsky-etal-2024-steering} and are less compute and data intensive than past work in readability controlled summarization~\cite{ribeiro-etal-2023-generating}. %\todoit{despite the attractiveness, we need to deal iwth previous} 
However, despite their appeal, prior work has noted challenges in reliability of these vectors ~\parencite{vonrutte2024languagemodelsguidelatent, stickland2024steeringeffectsimprovingpostdeployment, 10.5555/3737916.3742333, braun2025understandingunreliabilitysteeringvectors}. In this chapter, we propose a requirements based system for data selection that improves on these failure points.

% \todoit{make our task more clear}
% \todoit{make it more clear that this is an insight - in order to get these to work in our setting, we show that you need X} 


Our goal is to apply control vectors to the task of readability controlled summarization. We hypothesize that, in order for them to be effective, the extraction data needs to meet two Requirements (Req.): (1) separability of the paired examples and (2) paired examples contain comparable content. Regarding Req. 1, if the paired examples are too similar in readability, the control vectors will not have enough information, and will therefore be ineffective. Regarding Req. 2, control vectors work by taking the difference between the hidden states of the positive and negative examples. If the content of the summaries differs, the resulting vector will encode factors beyond readability, reducing its effectiveness. 
% \todoit{we don't want to say - this method doesn't work on most data. instead - for this to work, we need X - it turns out most datasets don't meet these requirements and these requirements would affect other methods}
Under these simple requirements, control vectors extracted from compliant data can interpolate between plain and technical summaries and generalize to other datasets.

We conduct an analysis of PLS datasets, and find that many popular datasets do not satisfy our proposed requirements. While we show that non-compliant data yields poor control vectors, such issues likely affect other summarization methods as well. Because control vectors are a low data method, applying our approach to a new task requires collecting a small curated dataset.

%\toditit{be careful here - say something like the advantages of our methods} 
Our work clarifies the data requirements for control vectors, increasing the chance of future success for new tasks. 